\documentclass[11pt]{article}
% ===========================================================================
%  Shared preamble for the AP Calculus AB exam-preparation worksheets.
%  Input right after \documentclass. Files beginning with "_" are not compiled
%  on their own. Navy and gold, to match the course's lesson pages.
% ===========================================================================
\usepackage[letterpaper,top=2.3cm,bottom=2.2cm,left=1.7cm,right=1.7cm,headheight=28pt,footskip=30pt]{geometry}
\usepackage{amsmath,amssymb}
\usepackage[T1]{fontenc}
\usepackage[utf8]{inputenc}
\usepackage{lmodern}
\usepackage{xcolor}
\usepackage[most]{tcolorbox}
\usepackage{enumitem}
\usepackage{fancyhdr}
\usepackage{array}
\usepackage{tikz}
\usepackage{pgfplots}
\pgfplotsset{compat=1.18}

\definecolor{navy}{HTML}{161E54}
\definecolor{ink}{HTML}{27357E}
\definecolor{gold}{HTML}{B98A00}
\definecolor{goldbg}{HTML}{FFFDF3}
\definecolor{inkbg}{HTML}{F6F7FD}
\definecolor{rule}{HTML}{C9CFF0}
\definecolor{okgreen}{HTML}{1B7A44}
\definecolor{okbg}{HTML}{F1F8F3}

\pagestyle{fancy}
\fancyhf{}
\fancyhead[L]{\footnotesize Name: \underline{\hspace{4cm}}}
\fancyhead[C]{\textbf{\color{navy}Integration Academy}\\[-2pt]{\scriptsize Calculus AB \textperiodcentered{} AP\textsuperscript{\textregistered} Exam Preparation}}
\fancyfoot[L]{\tiny\color{black!55}\parbox[t]{0.78\textwidth}{AP\textsuperscript{\textregistered} is a registered trademark of the College Board, which was not involved in the production of, and does not endorse, this material. \textcopyright{} Integration Academy.}}
\fancyfoot[R]{\scriptsize\thepage}
\renewcommand{\headrulewidth}{1.2pt}
\renewcommand{\headrule}{\hbox to\headwidth{\color{gold}\leaders\hrule height \headrulewidth\hfill}}
\renewcommand{\footrulewidth}{0pt}
\setlength{\parindent}{0pt}
\setlength{\parskip}{4pt}

% Title block: \wstitle{code}{title}{unit line}
\newcommand{\wstitle}[3]{\begin{center}{\small\bfseries\color{gold}WORKSHEET #1}\\[3pt]{\LARGE\bfseries\color{navy}#2}\\[3pt]{\small\color{black!60}#3}\end{center}\vspace{2pt}}
% Section heading
\newcommand{\sech}[1]{\par\vspace{10pt}{\large\bfseries\color{navy}#1}\par\nobreak\vspace{2pt}{\color{rule}\hrule height0.8pt}\vspace{7pt}}
% A level tag printed beside a group of questions
\newcommand{\level}[1]{\par\vspace{6pt}{\footnotesize\bfseries\color{gold}\MakeUppercase{#1}}\par\nobreak\vspace{2pt}}

\newtcolorbox{ideabox}{colframe=ink,colback=inkbg,coltitle=white,colbacktitle=ink,fonttitle=\bfseries\small,title={Key ideas},arc=2pt,boxrule=0.8pt,left=8pt,right=8pt}
\newtcolorbox{example}[1]{colframe=rule,colback=white,coltitle=navy,colbacktitle=inkbg,fonttitle=\bfseries\small,title={#1},arc=2pt,boxrule=0.8pt,borderline west={3pt}{0pt}{ink},left=10pt,right=8pt}
\newtcolorbox{frbox}[1]{colframe=gold,colback=goldbg,coltitle=white,colbacktitle=gold,fonttitle=\bfseries\small,title={#1},arc=2pt,boxrule=0.8pt,breakable,left=8pt,right=8pt}
\newtcolorbox{solbox}{colframe=okgreen,colback=okbg,coltitle=white,colbacktitle=okgreen,fonttitle=\bfseries\small,title={Worked solutions},arc=2pt,boxrule=0.8pt,breakable,left=8pt,right=8pt}
\newtcolorbox{scorebox}[1]{colframe=navy,colback=white,coltitle=white,colbacktitle=navy,fonttitle=\bfseries\small,title={#1},arc=2pt,boxrule=0.8pt,breakable,left=8pt,right=8pt}

\newcommand{\soln}{\par\smallskip\textbf{\color{ink}Solution.}\ }
\newcommand{\result}[1]{\par\smallskip\textbf{#1}}
% blank working space of a given height
\newcommand{\work}[1]{\par\vspace{#1}}
% a point value in the margin of a scoring guide
\newcommand{\pts}[1]{\hfill{\footnotesize\bfseries\color{gold}[#1]}}

% Coordinate grid for graph questions: \begin{apgraph}{xmin}{xmax}{ymin}{ymax} ... \end{apgraph}
\newenvironment{apgraph}[4]{\begin{center}\begin{tikzpicture}\begin{axis}[
  axis lines=middle, xmin=#1, xmax=#2, ymin=#3, ymax=#4,
  grid=both, grid style={black!12}, xtick distance=1, ytick distance=1,
  ticklabel style={font=\scriptsize}, xlabel={$x$}, ylabel={$y$},
  width=10cm, height=7.2cm, samples=80, clip=false,
  every axis plot/.append style={line width=1.1pt, color=ink}]}{\end{axis}\end{tikzpicture}\end{center}}
\newcommand{\closeddot}[2]{\addplot[only marks, mark=*, mark size=2.2pt, color=ink] coordinates {(#1,#2)};}
\newcommand{\opendot}[2]{\addplot[only marks, mark=*, mark size=2.2pt, color=ink, mark options={fill=white}] coordinates {(#1,#2)};}

% Multiple-choice options: four across, or stacked when the options are long.
\newcommand{\opts}[4]{\par\vspace{3pt}\begin{tabular}{@{}p{0.235\linewidth}p{0.235\linewidth}p{0.235\linewidth}p{0.235\linewidth}@{}}(A)\ #1 & (B)\ #2 & (C)\ #3 & (D)\ #4\end{tabular}\par\vspace{9pt}}
\newcommand{\optsv}[4]{\par\vspace{3pt}\begin{tabular}{@{}l@{\ \ }p{0.9\linewidth}@{}}(A) & #1\\(B) & #2\\(C) & #3\\(D) & #4\end{tabular}\par\vspace{9pt}}

\fancyhead[R]{\footnotesize\color{black!60}Worksheet 6.5}
\begin{document}

\wstitle{6.5}{Analysing a Function Defined by an Integral}{Unit 6: Integration and Accumulation of Change}

\begin{ideabox}
\begin{itemize}[leftmargin=*,itemsep=2pt,topsep=2pt]
  \item If $g(x)=\displaystyle\int_{a}^{x}f(t)\,dt$, then $g'(x)=f(x)$ and $g''(x)=f'(x)$.
  \item $g$ is increasing where $f$ is positive; $g$ has a relative maximum where $f$ changes from positive to negative.
  \item The graph of $g$ is concave up where $f$ is increasing, and has a point of inflection where $f$ changes direction.
  \item Values of $g$ are signed areas under the graph of $f$.
\end{itemize}
\end{ideabox}

\sech{Worked examples}
\begin{example}{Example 1 \textperiodcentered{} A relative minimum}
Let $g(x)=\displaystyle\int_{0}^{x}(t-3)\,dt$. Find the relative minimum of $g$.\soln $g'(x)=x-3$ changes from negative to positive at $x=3$. $g(3)$ is the signed area of a triangle below the axis.\result{Relative minimum $g(3)=-\dfrac92$.}
\end{example}

\begin{example}{Example 2 \textperiodcentered{} Concavity}
If $g'=f$ and $f$ is decreasing on $(2,5)$, what is true of the graph of $g$ there?\soln $g''=f'<0$.\result{It is concave down.}
\end{example}

\begin{example}{Example 3 \textperiodcentered{} An absolute extremum}
Let $g(x)=\displaystyle\int_{0}^{x}\left(4-t^2\right)\,dt$. Find the absolute maximum of $g$ on $[0,3]$.\soln $g'(x)=4-x^2=0$ at $x=2$. $g(0)=0$, $g(2)=8-\dfrac83=\dfrac{16}{3}$, $g(3)=12-9=3$.\result{$\dfrac{16}{3}$, at $x=2$.}
\end{example}

\newpage
\sech{Practice}
Questions 1 to 8 use this graph of $f$, which consists of three line segments. Let $g(x)=\displaystyle\int_{0}^{x}f(t)\,dt$.
\begin{apgraph}{-0.5}{9}{-3}{3}
\addplot[color=ink, mark=none, line width=1.1pt] coordinates {(0,-2) (2,2) (6,-2) (8,0)};
\closeddot{0}{-2}\closeddot{8}{0}
\node[font=\small, color=ink] at (axis cs:3.3,1.9) {$y=f(t)$};
\end{apgraph}
\level{Core}
\begin{enumerate}[leftmargin=*,itemsep=2pt]
  \item Find $g(2)$.\work{1.1cm}
  \item Find $g'(2)$.\work{1.1cm}
  \item Find $g(4)$.\work{1.1cm}
  \item Find $g''(3)$.\work{1.1cm}
\end{enumerate}

\level{Exam level}
\begin{enumerate}[leftmargin=*,itemsep=2pt,start=5]
  \item Find the $x$-coordinates of the relative minimum and the relative maximum of $g$ on $(0,8)$. Justify your answers.\work{2.1cm}
  \item Find the absolute maximum and absolute minimum values of $g$ on $[0,8]$.\work{2.1cm}
  \item On what intervals is the graph of $g$ concave up? Where are its points of inflection?\work{2.1cm}
  \item Write an equation of the line tangent to the graph of $g$ at $x=6$.\work{2.1cm}
\end{enumerate}

\begin{frbox}{Question 9 \textperiodcentered{} Free-response practice}
Let $H(x)=\displaystyle\int_{1}^{x}\ln t\,dt$ for $x>0$.
\begin{enumerate}[label=(\alph*),leftmargin=*,itemsep=2pt]
  \item Find $H(1)$ and $H'(x)$.\work{1.9cm}
  \item Show that $H$ has its absolute minimum at $x=1$.\work{1.9cm}
  \item Describe the concavity of the graph of $H$. Given that $H(e)=1$, write an equation of the tangent line at $x=e$.\work{1.9cm}
\end{enumerate}
\end{frbox}

\newpage
\begin{solbox}
\begin{enumerate}[leftmargin=*,itemsep=4pt]
  \item The areas on $[0,1]$ and $[1,2]$ are $-1$ and $1$: $g(2)=0$.
  \item $g'(2)=f(2)=2$.
  \item $g(2)+2=2$ (a triangle of area 2 on $[2,4]$).
  \item The slope of $f$ on $(2,6)$: $-1$.
  \item Minimum at $x=1$, where $g'=f$ changes from negative to positive. Maximum at $x=4$, where $f$ changes from positive to negative.
  \item $g(0)=0$, $g(1)=-1$, $g(4)=2$, $g(8)=2-2-2=-2$. Maximum 2; minimum $-2$ (at $x=8$).
  \item Concave up on $(0,2)$ and $(6,8)$, where $f$ is increasing. Points of inflection at $x=2$ and $x=6$.
  \item $g(6)=0$ and $g'(6)=f(6)=-2$: $y=-2(x-6)$.
  \item \begin{enumerate}[label=(\alph*),leftmargin=*,itemsep=3pt]
    \item $H(1)=0$ and $H'(x)=\ln x$.
    \item $H'(x)<0$ for $0<x<1$ and $H'(x)>0$ for $x>1$, so $H$ decreases and then increases: the absolute minimum is $H(1)=0$.
    \item $H''(x)=\dfrac1x>0$: concave up for all $x>0$. The slope at $x=e$ is $\ln e=1$: $y-1=x-e$.
  \end{enumerate}
\end{enumerate}
\end{solbox}

\end{document}
